\section{Bases de Diseño: Contexto del Sitio y del Sistema}

\subsection{Condiciones Ambientales del Sitio}

Las condiciones ambientales de diseño del sitio, Cajicá (Cundinamarca), se presentan en la Tabla \ref{tab:inf002_condiciones}. La altitud del casco urbano se validó contra el aeródromo Guaymaral (SKGY, 8 389 ft, a aproximadamente 7 km del sitio) \cite{dbcity_cajica} y la presión atmosférica se calculó con el modelo de atmósfera estándar ISA a 2 558 msnm, validada contra el QNH del METAR de El Dorado (SKBO) \cite{flightaware_skbo}.

\begin{table}[htbp]
	\centering
	\caption{Condiciones ambientales de diseño del sitio (consulta 23/07/2026)}
	\label{tab:inf002_condiciones}
	\setcounter{itemcount}{0}
	\small
	\begin{tabularx}{\textwidth}{@{}NXlc@{}}
		\toprule
		\multicolumn{1}{c}{\textbf{Item}} & \textbf{Parámetro} & \textbf{Valor de Diseño} & \textbf{Fuente} \\
		\midrule
		& Altitud (casco urbano Cajicá) & 2 558 msnm & \cite{dbcity_cajica} \\
		& Presión atmosférica & 74.1 kPa & ISA a 2 558 m \cite{flightaware_skbo} \\
		& Temperatura de diseño verano & 21 $^{\circ}$C (condición 0.4 \%, bulbo seco) & \cite{captiveaire_ashrae2009} \\
		& Temperatura de diseño invierno & 3 $^{\circ}$C (condición 99.6 \%) & \cite{captiveaire_ashrae2009} \\
		& Temperatura media anual & 14 $^{\circ}$C & \cite{alcaldia_cajica} \\
		& Humedad relativa media anual & 84 \% (rango mensual 80--89 \%) & \cite{weatheratlas_cajica} \\
		& Punto de rocío de diseño & 14.2 $^{\circ}$C (deshumidificación 0.4 \%) & \cite{captiveaire_ashrae2009} \\
		& Densidad del aire & 0.88 kg/m$^3$ a 20 $^{\circ}$C; 0.87 kg/m$^3$ en condición 0.4 \% & Gas ideal, $P_{atm}$ = 74.1 kPa \\
		\bottomrule
	\end{tabularx}
\end{table}

La presión de estación real es 74.1 kPa; los valores de 1 024--1 029 hPa publicados en portales meteorológicos corresponden a presión reducida a nivel del mar (QNH) y no deben utilizarse en el diseño, pues confundir ambas introduce un error del 28 \% en la densidad del aire y en la selección del ventilador \cite{tutiempo_cajica}. Los valores ASHRAE citados corresponden a la edición 2009, única tabla extraíble de la fuente; la edición 2021/2025 puede diferir en $\pm$0.5 $^{\circ}$C, diferencia irrelevante para un sistema de ventilación sin climatización. La estación proxy (El Dorado, 2 546 msnm) se encuentra prácticamente a la cota del sitio.

\subsection{Corrección de Densidad y su Impacto en la Selección}

Las curvas de catálogo de ventiladores se publican a densidad estándar (1.2 kg/m$^3$; 70 $^{\circ}$F, 29.92 in Hg), condición declarada explícitamente por los fabricantes \cite{tcf_fe1600}. A tamaño y velocidad constantes, el caudal volumétrico no cambia con la densidad, mientras que la presión y la potencia absorbida son proporcionales a ella. El factor de densidad del sitio es:

\begin{equation}
	k = \frac{\rho_{sitio}}{\rho_{est}} = \frac{0.88}{1.20} = 0.733
	\label{eq:inf002_kdensidad}
\end{equation}

En consecuencia, todo equipo seleccionado por catálogo a densidad estándar entregará en el sitio la misma presión estática reducida en ese factor, y las caídas de presión de filtros, rejillas y accesorios tabuladas a densidad estándar deben multiplicarse por 0.733 para obtener el valor en el sitio. La Tabla \ref{tab:inf002_punto_trabajo} cuantifica el impacto sobre el punto de trabajo, con los valores congelados de la memoria de cálculo.

\begin{table}[htbp]
	\centering
	\caption{Punto de trabajo del sistema en el sitio y equivalente de catálogo}
	\label{tab:inf002_punto_trabajo}
	\setcounter{itemcount}{0}
	\small
	\begin{tabularx}{\textwidth}{@{}NXcc@{}}
		\toprule
		\multicolumn{1}{c}{\textbf{Item}} & \textbf{Magnitud} & \textbf{En el Sitio} & \textbf{Equivalente de Catálogo} \\
		& & ($\rho$ = 0.88 kg/m$^3$) & ($\rho$ = 1.2 kg/m$^3$) \\
		\midrule
		& Caudal & 3 840 m$^3$/h (1.0667 m$^3$/s; 2 260 CFM) & 3 840 m$^3$/h (sin cambio) \\
		& $\Delta P$ filtro MERV 13-14 cargado (diseño) & 154 Pa & 210 Pa \\
		& $\Delta P$ filtro MERV 13-14 limpio & 59 Pa & 80 Pa \\
		& Offset de presurización & 25 Pa & 34 Pa \\
		& $\Delta P$ rejillas (3 m/s, $C_d$ = 0.60) & 11 Pa & 15 Pa \\
		& \textbf{$\Delta P$ total escenario cargado (selección)} & \textbf{190 Pa} & \textbf{260 Pa} \\
		& $\Delta P$ total escenario limpio & 95 Pa & 130 Pa \\
		& Potencia teórica de aire ($\eta$ = 0.60) & 0.338 kW (0.45 HP) & -- \\
		\bottomrule
	\end{tabularx}
\end{table}

La corrección tiene una consecuencia directa sobre la exfiltración: la configuración alternativa sin damper (rejillas a 4 m/s) ya no alcanza los 25 Pa de consigna a esta altitud, pues solo sostiene 19.6 Pa; para cerrar el balance en 25 Pa sin damper se requeriría un área total de 0.236 m$^2$ (0.079 m$^2$ por rejilla, aproximadamente 280 mm $\times$ 281 mm, a 4.5 m/s). En consecuencia, el damper de alivio barométrico pasa de opción a componente obligatorio del sistema.

Sobre el motor, la potencia teórica corregida al sitio es 0.338 kW (0.45 HP); se instala un motor de 1.0 HP TEFC con tratamiento anticorrosivo, lo que deja un margen superior al 100 \% sobre el eje, suficiente para absorber el derateo por altitud (regla práctica NEMA MG-1: 3 \% por cada 500 m sobre 1 000 msnm, es decir, aproximadamente 9 \% a 2 558 msnm \cite{nema_mg1_derate,silvertip_altitude}) y el factor de servicio exigido por el ambiente corrosivo.

\subsection{Corrosividad del Ambiente Exterior}

El laboratorio se ubica dentro de una planta de hipoclorito de calcio; la atmósfera exterior combina \ce{Cl2} y \ce{ClO-} (oxidantes fuertes), polvo de \ce{Ca(ClO)2} y humedad alta (84 \% media). La jerarquía de materiales documentada para este servicio es: PRFV con resina viniléster como primera opción (las resinas Derakane se especifican expresamente para cloro e hipoclorito \cite{ineos_derakane}), polipropileno como alternativa sólida a temperatura ambiente, acero con recubrimiento epóxico solo como opción presupuestal con inspección programada, y acero galvanizado descartado. El acero inoxidable 304/316 en la corriente de cloro sufre picadura y agrietamiento por corrosión bajo tensión, por lo que no es primera opción para la carcasa del ventilador, aunque se acepta inox 316L en elementos de baja solicitación química directa (rejillas, mallas, ferretería) por disponibilidad y fabricación local \cite{systech_corrosive}.

La regla de especificación adoptada para todo el sistema se resume en la Tabla \ref{tab:inf002_materiales}.

\begin{table}[htbp]
	\centering
	\caption{Regla de especificación de materiales del sistema}
	\label{tab:inf002_materiales}
	\setcounter{itemcount}{0}
	\small
	\begin{tabularx}{\textwidth}{@{}NlX@{}}
		\toprule
		\multicolumn{1}{c}{\textbf{Item}} & \textbf{Elemento} & \textbf{Especificación de Material} \\
		\midrule
		& Máquina de impulsión & PRFV viniléster o PP \\
		& Marcos de filtro y portafiltros & Marcos plásticos; portafiltros inox 304/316 \\
		& Rejillas y malla anti-insectos & Inox 316L \\
		& Damper de alivio & Aluminio extruido 6063-T5 o línea severe-environment \\
		& Ferretería y soportes & Inox 316 (A4) con aislamiento dieléctrico frente a estructura galvanizada \\
		& Sellantes & Silicona RTV neutra \\
		& Conexión flexible & Hipalón (CSM) \\
		& Instrumentos hacia el exterior & Gabinete IP66 \\
		\bottomrule
	\end{tabularx}
\end{table}
